\section{Formulario}

\subsection{Trabajo de la fuerza eléctrica}

Para un desplazamiento desde \(A\) hasta \(B\),

\[
\boxed{
W_{e,A\to B}
=
\int_A^B \vec F_e\cdot d\vec r
}
\]

En componentes rectangulares,

\[
\boxed{
W_{e,A\to B}
=
\int_A^B
\left(
F_x\,dx+F_y\,dy+F_z\,dz
\right)
}
\]

Como

\[
F_x=qE_x,\qquad
F_y=qE_y,\qquad
F_z=qE_z,
\]

también puede escribirse

\[
\boxed{
W_{e,A\to B}
=
q\int_A^B
\left(
E_x\,dx+E_y\,dy+E_z\,dz
\right)
}.
\]


\subsection{Fuerza electrostática conservativa}

El trabajo entre dos puntos es independiente de la trayectoria.

Para una trayectoria cerrada,

\[
\boxed{
\oint\vec F_e\cdot d\vec r=0
}.
\]


\subsection{Cambio de energía potencial eléctrica}

\[
\boxed{
\Delta U=U_f-U_i
}
\]

La relación con el trabajo de la fuerza eléctrica es

\[
\boxed{
\Delta U=-W_e
}
\]

o

\[
\boxed{
W_e=U_i-U_f
}.
\]

Por tanto,

\[
W_e>0
\Longrightarrow
\Delta U<0,
\]

\[
W_e<0
\Longrightarrow
\Delta U>0.
\]

Unidad SI:

\[
\boxed{[U]=\si{\joule}}.
\]


\subsection{Trabajo externo}

Si una fuerza externa desplaza lentamente las cargas y no cambia apreciablemente
su energía cinética,

\[
\boxed{
W_{\text{ext}}=\Delta U
}
\]

y

\[
\boxed{
W_{\text{ext}}=-W_e
}.
\]


\subsection{Cambio de energía potencial a partir del campo}

En general,

\[
\boxed{
\Delta U
=
-q\int_A^B
\left(
E_x\,dx+E_y\,dy+E_z\,dz
\right)
}.
\]

Para un campo uniforme,

\[
\boxed{
\Delta U
=
-q
\left(
E_x\Delta x+
E_y\Delta y+
E_z\Delta z
\right)
}.
\]

En una dimensión,

\[
\boxed{
\Delta U=-qE_x\Delta x
}.
\]

El trabajo eléctrico correspondiente es

\[
\boxed{
W_e=
q
\left(
E_x\Delta x+
E_y\Delta y+
E_z\Delta z
\right)
}.
\]


\subsection{Referencia de energía potencial}

Para sistemas de cargas puntuales aisladas se utiliza normalmente

\[
\boxed{
U(\infty)=0
}.
\]


\subsection{Energía potencial de dos cargas puntuales}

Para dos cargas \(q_1\) y \(q_2\) separadas una distancia \(r\),

\[
\boxed{
U=k\frac{q_1q_2}{r}
}
\]

con

\[
\boxed{
k\approx
\SI{8.99e9}{\newton\metre\squared\per\coulomb\squared}
}.
\]

También,

\[
\boxed{
k=\frac{1}{4\pi\varepsilon_0}
}.
\]

En esta expresión no se utiliza valor absoluto:

\[
\boxed{
U=k\frac{q_1q_2}{r}
}.
\]


\subsection{Signo de la energía potencial}

Cargas del mismo signo:

\[
q_1q_2>0
\Longrightarrow
\boxed{U>0}.
\]

Cargas de signos opuestos:

\[
q_1q_2<0
\Longrightarrow
\boxed{U<0}.
\]

Cuando

\[
r\to\infty,
\]

\[
\boxed{U\to0}.
\]


\subsection{Dependencia con la distancia}

Para dos cargas puntuales,

\[
\boxed{
U\propto\frac{1}{r}
}.
\]

Para cargas del mismo signo:

\[
r\downarrow
\Longrightarrow
U\uparrow.
\]

Para cargas de signos opuestos:

\[
r\downarrow
\Longrightarrow
U\text{ se hace más negativa}.
\]


\subsection{Cambio de separación entre dos cargas}

Si la separación cambia de \(r_i\) a \(r_f\),

\[
U_i=k\frac{q_1q_2}{r_i},
\qquad
U_f=k\frac{q_1q_2}{r_f}.
\]

Entonces,

\[
\boxed{
\Delta U
=
kq_1q_2
\left(
\frac{1}{r_f}
-
\frac{1}{r_i}
\right)
}.
\]

El trabajo realizado por la fuerza eléctrica es

\[
\boxed{
W_e
=
kq_1q_2
\left(
\frac{1}{r_i}
-
\frac{1}{r_f}
\right)
}.
\]

El trabajo externo cuasiestático es

\[
\boxed{
W_{\text{ext}}
=
kq_1q_2
\left(
\frac{1}{r_f}
-
\frac{1}{r_i}
\right)
}.
\]


\subsection{Ensamblaje desde el infinito}

Si inicialmente

\[
r_i\to\infty,
\]

entonces

\[
U_i=0.
\]

Para llevar lentamente las cargas hasta una separación \(r\),

\[
\boxed{
W_{\text{ext}}
=
U
=
k\frac{q_1q_2}{r}
}.
\]

El trabajo realizado por la fuerza eléctrica es

\[
\boxed{
W_e
=
-k\frac{q_1q_2}{r}
}.
\]


\subsection{Separación de cargas de signos opuestos}

Si dos cargas de signos opuestos están inicialmente separadas una distancia \(r\)
y se llevan lentamente hasta el infinito,

\[
U_i=k\frac{q_1q_2}{r},
\qquad
U_f=0.
\]

Entonces,

\[
\boxed{
W_{\text{ext}}
=
-k\frac{q_1q_2}{r}
}.
\]

Como

\[
q_1q_2<0,
\]

se cumple

\[
W_{\text{ext}}>0.
\]


\subsection{Energía potencial de varias cargas}

Para \(N\) cargas puntuales,

\[
\boxed{
U
=
k\sum_{i<j}
\frac{q_iq_j}{r_{ij}}
}.
\]

Cada par debe contarse una sola vez.


\subsection{Sistema de tres cargas}

\[
\boxed{
U
=
k
\left(
\frac{q_1q_2}{r_{12}}
+
\frac{q_1q_3}{r_{13}}
+
\frac{q_2q_3}{r_{23}}
\right)
}.
\]


\subsection{Sistema de cuatro cargas}

Para cuatro cargas existen seis interacciones:

\[
\boxed{
U
=
k
\left(
\frac{q_1q_2}{r_{12}}
+
\frac{q_1q_3}{r_{13}}
+
\frac{q_1q_4}{r_{14}}
+
\frac{q_2q_3}{r_{23}}
+
\frac{q_2q_4}{r_{24}}
+
\frac{q_3q_4}{r_{34}}
\right)
}.
\]


\subsection{Número de pares}

Para \(N\) cargas,

\[
\boxed{
N_{\text{pares}}
=
\frac{N(N-1)}{2}
}.
\]


\subsection{Distancias entre cargas}

En una dimensión,

\[
\boxed{
r_{ij}=|x_j-x_i|
}.
\]

En dos dimensiones,

\[
\boxed{
r_{ij}
=
\sqrt{
(x_j-x_i)^2+
(y_j-y_i)^2
}
}.
\]

En tres dimensiones,

\[
\boxed{
r_{ij}
=
\sqrt{
(x_j-x_i)^2+
(y_j-y_i)^2+
(z_j-z_i)^2
}
}.
\]


\subsection{Conservación de la energía}

Si sólo intervienen fuerzas conservativas electrostáticas,

\[
\boxed{
K_i+U_i=K_f+U_f
}.
\]

Equivalentemente,

\[
\boxed{
\Delta K+\Delta U=0
}
\]

o

\[
\boxed{
\Delta K=-\Delta U
}.
\]


\subsection{Energía cinética}

Para una partícula de masa \(m\) y rapidez \(v\),

\[
\boxed{
K=\frac12mv^2
}.
\]


\subsection{Partícula liberada desde el reposo}

Si

\[
K_i=0,
\]

entonces

\[
U_i=K_f+U_f.
\]

Por tanto,

\[
\boxed{
K_f=U_i-U_f
}.
\]

Como

\[
K_f=\frac12mv_f^2,
\]

resulta

\[
\boxed{
\frac12mv_f^2=U_i-U_f
}
\]

y

\[
\boxed{
v_f=
\sqrt{
\frac{2(U_i-U_f)}{m}
}
}.
\]


\subsection{Rapidez entre dos posiciones}

Si la partícula posee inicialmente rapidez \(v_i\),

\[
\frac12mv_i^2+U_i
=
\frac12mv_f^2+U_f.
\]

Despejando,

\[
\boxed{
v_f
=
\sqrt{
v_i^2+
\frac{2(U_i-U_f)}{m}
}
}.
\]


\subsection{Dos cargas y conservación de la energía}

Para una carga móvil \(q\) que interactúa con una carga fija \(Q\),

\[
U_i=k\frac{Qq}{r_i},
\qquad
U_f=k\frac{Qq}{r_f}.
\]

Entonces,

\[
\boxed{
\frac12mv_i^2
+
k\frac{Qq}{r_i}
=
\frac12mv_f^2
+
k\frac{Qq}{r_f}
}.
\]

Si parte del reposo,

\[
\boxed{
\frac12mv_f^2
=
kQq
\left(
\frac{1}{r_i}
-
\frac{1}{r_f}
\right)
}.
\]


\subsection{Movimiento espontáneo}

Si sólo actúa la interacción eléctrica,

\[
\boxed{
U\downarrow
\Longrightarrow
K\uparrow
}.
\]

Para cargas del mismo signo,

\[
\boxed{
r\uparrow,\qquad U\downarrow,\qquad K\uparrow
}.
\]

Para cargas de signos opuestos,

\[
\boxed{
r\downarrow,\qquad U\downarrow,\qquad K\uparrow
}.
\]


\subsection{Fuerza a partir de la energía potencial}

Para una interacción radial,

\[
\boxed{
F_r=-\frac{dU}{dr}
}.
\]

Si

\[
U=k\frac{q_1q_2}{r},
\]

entonces

\[
\boxed{
F_r=k\frac{q_1q_2}{r^2}
}.
\]

El signo indica el sentido radial:

\[
q_1q_2>0
\Longrightarrow
F_r>0
\]

corresponde a repulsión,

\[
q_1q_2<0
\Longrightarrow
F_r<0
\]

corresponde a atracción.


\subsection{Procedimiento para dos cargas}

\begin{enumerate}
    \item Identificar \(q_1\), \(q_2\) y la separación \(r\).
    \item Convertir las magnitudes al SI.
    \item Conservar los signos de las cargas.
    \item Aplicar
    \[
    U=k\frac{q_1q_2}{r}.
    \]
    \item Interpretar el signo del resultado.
\end{enumerate}


\subsection{Procedimiento para un sistema de varias cargas}

\begin{enumerate}
    \item Identificar todas las cargas.
    \item Determinar las distancias entre cada par.
    \item Calcular
    \[
    U_{ij}=k\frac{q_iq_j}{r_{ij}}.
    \]
    \item Conservar el signo de cada producto \(q_iq_j\).
    \item Contar cada par una sola vez.
    \item Sumar:
    \[
    U=\sum_{i<j}U_{ij}.
    \]
    \item Interpretar el signo y el valor de la energía total.
\end{enumerate}


\subsection{Procedimiento para problemas de movimiento}

\begin{enumerate}
    \item Definir el estado inicial y el estado final.
    \item Calcular \(U_i\).
    \item Calcular \(U_f\).
    \item Determinar \(K_i\).
    \item Aplicar
    \[
    K_i+U_i=K_f+U_f.
    \]
    \item Despejar la rapidez o la magnitud solicitada.
    \item Comprobar que la energía cinética obtenida sea físicamente posible:
    \[
    K_f\geq0.
    \]
\end{enumerate}


\subsection{Advertencias}

\begin{itemize}
    \item No utilizar valor absoluto en
    \[
    U=k\frac{q_1q_2}{r}.
    \]

    \item La energía potencial es escalar; no tiene dirección.

    \item Una energía potencial negativa es físicamente válida.

    \item El valor de \(U\) depende de la referencia elegida.

    \item Para cargas puntuales aisladas suele utilizarse
    \[
    U(\infty)=0.
    \]

    \item El trabajo realizado por la fuerza eléctrica satisface
    \[
    W_e=-\Delta U.
    \]

    \item El trabajo externo cuasiestático satisface
    \[
    W_{\text{ext}}=\Delta U.
    \]

    \item En un sistema de varias cargas, cada interacción se cuenta una sola vez.

    \item Una carga neta total igual a cero no implica
    \[
    U=0.
    \]

    \item La conservación
    \[
    K_i+U_i=K_f+U_f
    \]
    requiere que no existan transferencias adicionales de energía que hayan sido
    omitidas.
\end{itemize}


\subsection{Resumen de fórmulas}

\begin{align*}
W_e
&=
\int_A^B
\left(
F_x\,dx+F_y\,dy+F_z\,dz
\right),\\
W_e
&=
q\int_A^B
\left(
E_x\,dx+E_y\,dy+E_z\,dz
\right),\\
\Delta U
&=-W_e,\\
W_{\text{ext}}
&=\Delta U,\\
\Delta U
&=
-q\int_A^B
\left(
E_x\,dx+E_y\,dy+E_z\,dz
\right),\\
\Delta U_{\text{uniforme}}
&=
-q(E_x\Delta x+E_y\Delta y+E_z\Delta z),\\
U
&=
k\frac{q_1q_2}{r},\\
\Delta U
&=
kq_1q_2
\left(
\frac{1}{r_f}-\frac{1}{r_i}
\right),\\
W_e
&=
kq_1q_2
\left(
\frac{1}{r_i}-\frac{1}{r_f}
\right),\\
U_{\text{sistema}}
&=
k\sum_{i<j}\frac{q_iq_j}{r_{ij}},\\
N_{\text{pares}}
&=
\frac{N(N-1)}{2},\\
K
&=
\frac12mv^2,\\
K_i+U_i
&=
K_f+U_f,\\
\Delta K
&=-\Delta U,\\
F_r
&=
-\frac{dU}{dr}.
\end{align*}